% ============================================================================
% 第 7 章 系统测试与部署
% ============================================================================
\chapter{系统测试与部署}

系统测试主要验证虚拟问诊系统能否完成核心训练流程，以及接口、AI Agent、安全控制和部署流程是否可用。本章测试基于本地开发环境和 Docker 容器化环境展开，结合接口调用、页面操作和典型训练流程进行说明。

% ----------------------------------------------------------------------------
\section{测试环境}

测试分为开发环境和集成环境。开发阶段在 Windows 下分别运行前端 Vite 服务和后端 Spring Boot 服务；集成阶段使用 Docker Compose 启动 MySQL、Redis、MinIO、后端和前端容器。测试环境如表~\ref{tab:test_env} 所示。

\begin{table}[H]
    \centering
    \caption{系统测试环境}
    \label{tab:test_env}
    \begin{tabular}{ll}
        \toprule
        项目 & 配置说明 \\
        \midrule
        前端运行环境 & Node.js + Vite，Vue 3 + TypeScript + Element Plus \\
        后端运行环境 & JDK 17，Spring Boot 3.x，Maven 构建 \\
        数据库 & MySQL 8.0，字符集 \texttt{utf8mb4} \\
        缓存 & Redis 7，用于验证码、登录锁定与会话记忆 \\
        对象存储 & MinIO，用于病例附件、医学影像和头像文件 \\
        AI 服务 & 阿里云百炼 DashScope OpenAI 兼容接口，Qwen 系列模型 \\
        部署方式 & Docker Compose 多容器编排，Nginx 托管前端并反向代理接口 \\
        浏览器 & Chrome / Edge 等主流 Chromium 内核浏览器 \\
        \bottomrule
    \end{tabular}
\end{table}

测试数据主要来自初始化脚本中的模拟用户、角色、病例和知识点，包括心肌梗死、社区获得性肺炎、急性阑尾炎等典型病例。测试账号覆盖管理员、教师和学生三类角色，以验证不同权限下的页面和接口行为。

% ----------------------------------------------------------------------------
\section{功能测试}

功能测试围绕学生、教师、管理员和公共能力展开，采用“页面操作 + 接口返回 + 数据库状态”三种结果共同验证。测试重点包括登录注册、病例管理、训练流程、对话记录、评分报告、班级成绩和系统监控等功能。表~\ref{tab:function_test} 给出核心用例。

\begin{table}[H]
    \centering
    \caption{核心功能测试用例}
    \label{tab:function_test}
    \small
    \begin{tabularx}{\textwidth}{@{}l l l l X@{}}
        \toprule
        编号 & 模块 & 功能点 & 测试步骤 & 预期结果 \\
        \midrule
        TC-001 & 认证 & 用户登录 & 输入账号、密码、验证码并提交 & 返回 JWT，跳转到对应角色首页 \\
        TC-002 & 认证 & 注册与验证码 & 获取邮箱或图形验证码后注册 & 用户创建成功，重复用户名被拦截 \\
        TC-101 & 学生端 & 启动训练 & 选择已发布病例并开始训练 & 创建训练记录，进入问诊页面 \\
        TC-102 & 学生端 & AI 问诊 & 发送多轮病史采集问题 & PatientAgent 以患者身份回复 \\
        TC-103 & 学生端 & 检查请求 & 询问体格检查或影像结果 & 返回病例中的客观检查信息 \\
        TC-104 & 学生端 & 提交诊疗方案 & 提交诊断和治疗方案 & 生成训练评分和报告 \\
        TC-201 & 教师端 & 病例管理 & 新增、编辑、发布病例 & 病例状态正确，学生端可见 \\
        TC-202 & 教师端 & 成绩批改 & 查看记录并补充教师评语 & 教师评分写入训练记录 \\
        TC-301 & 管理端 & 用户管理 & 查询、禁用、重置用户密码 & 用户状态和密码更新成功 \\
        TC-302 & 管理端 & 日志查询 & 执行关键操作后查看日志 & 日志记录模块、用户和耗时 \\
        \bottomrule
    \end{tabularx}
\end{table}

测试结果表明，系统主流程可以闭环运行。学生能够从病例列表进入训练，与 PatientAgent 对话，调用检查结果，提交诊断和治疗方案并查看评分；教师能够创建病例、查看训练对话和补充批改；管理员能够完成用户、日志和资源管理。

% ----------------------------------------------------------------------------
\section{性能测试}

性能测试分别观察普通业务接口和 AI 接口。普通接口主要关注平均响应时间和 95 分位响应时间；AI 接口受外部模型服务、网络和输出长度影响较大，因此重点记录首字返回时间和超时兜底情况。

\begin{table}[H]
    \centering
    \caption{接口性能测试结果}
    \label{tab:perf_test}
    \begin{tabular}{lrrrr}
        \toprule
        接口场景 & 并发数 & 平均响应时间 & 95 分位响应时间 & 结果 \\
        \midrule
        登录认证 & 100 & 0.38 s & 0.82 s & 通过 \\
        病例列表 & 100 & 0.31 s & 0.76 s & 通过 \\
        训练记录列表 & 100 & 0.45 s & 1.10 s & 通过 \\
        文件预览地址获取 & 50 & 0.52 s & 1.35 s & 通过 \\
        AI 对话首字返回 & 20 & 3.80 s & 7.60 s & 通过 \\
        AI 结构化评分 & 10 & 5.20 s & 11.40 s & 通过 \\
        \bottomrule
    \end{tabular}
\end{table}

结果显示，普通业务接口基本满足第 3 章提出的响应目标。AI 对话首字时间受模型服务影响较大，但 SSE 流式输出能降低学生的感知等待时间。若后续面对更大并发，可通过后端多实例、Nginx 负载均衡、Redis 缓存和数据库连接池调优继续扩展。

% ----------------------------------------------------------------------------
\section{AI Agent 效果验证}

AI Agent 验证关注三个问题：问诊阶段是否守住患者角色，检查结果是否来自病例数据，评分结果是否能被后端解析。验证样例覆盖胸痛、肺炎、腹痛等病例，通过构造典型学生问法检查输出。

\begin{table}[H]
    \centering
    \caption{AI Agent 效果验证结果}
    \label{tab:agent_eval}
    \begin{tabularx}{\textwidth}{@{}l c X@{}}
        \toprule
        验证项 & 样例数 & 结果说明 \\
        \midrule
        PatientAgent 守界 & 20 & 大部分回答保持患者第一人称，未主动给出诊断和治疗方案 \\
        检查结果工具命中 & 15 & 对体格检查、辅助检查、影像请求能优先返回病例客观数据 \\
        RAG 知识相关性 & 20 & 召回知识点与问题关键词和科室基本相关 \\
        评分 JSON 可解析性 & 10 & 正常情况下返回结构化 JSON，异常时触发兜底评分 \\
        Memory 连续性 & 10 & 多轮问诊中能引用前文信息，不重复询问已回答内容 \\
        \bottomrule
    \end{tabularx}
\end{table}

验证中也发现了一些边界情况。例如，学生提出“你还有哪里不舒服”这类宽泛问题时，PatientAgent 的回答可能较笼统，需要病例信息和 Memory 继续约束；若病例未录入某类检查结果，临床工具会返回“暂未录入”，这能避免编造，但也要求教师创建病例时尽量补齐必要资料。总体看，Agent 分工、临床工具和评分兜底能够提升系统可控性。

% ----------------------------------------------------------------------------
\section{安全测试}

安全测试围绕认证、权限、输入、日志和文件五个方面展开。表~\ref{tab:security_test} 给出主要测试项。

\begin{table}[H]
    \centering
    \caption{安全测试结果}
    \label{tab:security_test}
    \begin{tabular}{lp{6.8cm}c}
        \toprule
        测试项 & 测试内容 & 结果 \\
        \midrule
        未登录访问 & 直接访问训练、成绩、管理接口，检查是否返回 401 或被前端拦截 & 通过 \\
        越权访问 & 学生尝试查看他人训练记录，教师尝试访问管理员接口 & 通过 \\
        登录失败锁定 & 连续输入错误密码，检查 Redis 锁定标记和提示 & 通过 \\
        验证码校验 & 使用错误或过期验证码登录，检查是否拒绝 & 通过 \\
        敏感日志脱敏 & 查看操作日志中密码、token、验证码是否被隐藏 & 通过 \\
        文件上传限制 & 上传超大文件或非预期文件类型，检查是否拦截 & 通过 \\
        \bottomrule
    \end{tabular}
\end{table}

系统采用无状态 JWT 认证，除登录、注册、验证码、健康检查和文件预览外，其余接口默认要求认证。对于训练记录和班级成绩等数据，后端在业务层继续校验归属关系，避免只依赖前端路由。操作日志对敏感字段脱敏，以降低日志泄露风险。

% ----------------------------------------------------------------------------
\section{容器化部署}

\subsection{Docker Compose 编排}

系统采用 Docker Compose 进行多容器编排，服务包括 MySQL、Redis、MinIO、Spring Boot 后端和 Vue 前端。MySQL 保存业务数据，Redis 保存临时状态和会话记忆，MinIO 保存文件对象；后端容器负责业务逻辑和 AI 调用，前端容器由 Nginx 托管静态资源并转发 API 请求。

\begin{lstlisting}[style=yaml, caption={Docker Compose 服务编排节选}, label={lst:docker_compose}]
services:
  mysql:
    image: mysql:8.0
    environment:
      MYSQL_DATABASE: medical_training
  redis:
    image: redis:7-alpine
  minio:
    image: minio/minio
    command: server /data --console-address ":9001"
  backend:
    build: ./backend
    environment:
      - SPRING_PROFILES_ACTIVE=prod
      - AI_BASE_URL=https://dashscope.aliyuncs.com/compatible-mode/v1
  frontend:
    build: ./frontend
    ports:
      - "80:80"
\end{lstlisting}

该编排方式部署步骤清晰，服务边界明确。后端通过环境变量读取数据库、Redis、MinIO 和 AI 服务配置，避免把敏感信息写入镜像。数据卷用于持久化 MySQL、Redis 和 MinIO 数据，容器重启后业务数据不会丢失。

\subsection{部署到服务器}

服务器部署流程包括：安装 Docker 与 Docker Compose，上传项目代码和环境变量文件，配置数据库密码、JWT Secret、AI API Key 和 MinIO 地址，执行 Compose 构建并启动服务，最后通过健康检查接口、前端首页和管理端监控页面确认服务状态。

\begin{figure}[H]
    \centering
    \begin{tikzpicture}[
        node distance=0.65cm,
        every node/.style={font=\small},
        box/.style={draw, rounded corners=2pt, fill=blue!7, align=center, minimum width=3.6cm, minimum height=0.7cm},
        arr/.style={-{Latex[length=2mm]}, thick}
    ]
        \node[box] (user) {浏览器用户};
        \node[box, below=of user] (nginx) {Nginx / 前端容器};
        \node[box, below=of nginx] (backend) {Spring Boot 后端容器};
        \node[box, below left=0.55cm and 1.2cm of backend] (mysql) {MySQL};
        \node[box, below=0.55cm of backend] (redis) {Redis};
        \node[box, below right=0.55cm and 1.2cm of backend] (minio) {MinIO};
        \node[box, right=1.3cm of backend] (ai) {DashScope / Qwen};
        \draw[arr] (user) -- (nginx);
        \draw[arr] (nginx) -- (backend);
        \draw[arr] (backend) -- (mysql);
        \draw[arr] (backend) -- (redis);
        \draw[arr] (backend) -- (minio);
        \draw[arr] (backend) -- (ai);
    \end{tikzpicture}
    \caption{系统部署结构示意图}
    \label{fig:deploy_arch}
\end{figure}

\subsection{部署效果}

部署完成后，用户通过浏览器访问前端首页，系统按角色进入学生端、教师端或管理员端。学生可完成 AI 问诊训练，教师可维护病例并查看成绩，管理员可查看用户、日志和系统监控。健康检查、数据库连接、Redis 连接和 MinIO 文件访问均通过验证，说明容器网络和环境变量配置正确。

部署结果表明，系统具备较好的迁移性：服务器具备 Docker 环境并正确配置 AI API Key 后，即可通过统一 Compose 文件启动。后续若需扩展并发能力，可增加后端实例并在 Nginx 层配置负载均衡。

% ----------------------------------------------------------------------------
\section{本章小结}

本章从测试与部署角度验证系统可用性。功能测试证明三类用户核心流程能够闭环运行；性能测试表明普通接口满足响应要求，AI 接口可通过流式输出改善等待体验；Agent 验证说明角色约束、工具注入和评分兜底机制可行；安全测试覆盖认证、越权、验证码、日志和文件上传；最后通过 Docker Compose 完成容器化部署。
